
\minerva uses a finely segmented detector to record interactions of neutrinos produced by the NuMI beam line~\cite{Anderson:1998zza} at Fermilab.  
Data for this analysis come from $2.94\times 10^{20}$ protons on target taken between March 2010 and April 2012 
when the beam line produced a broadband neutrino beam peaked at \unit[3.5]{\GeV} with $>95\%$ \numu at the peak energy.
The \minerva detector is comprised of 120 hexagonal modules 
perpendicular to the $z$ axis, which is tilted \unit[58]{mrad} upwards with respect to the beam line~\cite{minerva_nim}.
There are four module types: active tracking, 
electromagnetic calorimeter, hadronic calorimeter, and inactive nuclear target.
The most upstream part of
the detector includes five inactive targets, numbered from upstream to downstream, 
each separated by four active tracking modules.
Target 4 is lead; other targets comprise two or three materials arranged at differing transverse positions filling the $x-y$ plane.
Targets 1, 2, and 5 are constructed of steel and lead plates joined together;
target 3 has graphite, steel, and lead plates.  
Total fiducial masses of C, Fe, and Pb in the nuclear target region 
are 0.159, 0.628, and 0.711 tons, respectively.
A fully active tracking region with a fiducial mass of 5.48 tons is downstream of the nuclear target region.
The target and tracker regions are surrounded by electromagnetic and hadronic calorimeters.
The MINOS near detector, a magnetized iron spectrometer~\cite{Michael:2008bc}, is located 
\unit[2]{m} downstream of the \minerva detector.

Neutrino flux is predicted using a GEANT4-based simulation tuned to hadron production data~\cite{Alt:2006fr} as described in Ref.~\cite{nubarprl}\SuppNote{See Supplemental Material\SuppLocation for a table of the simulated neutrino flux}.
Neutrino interactions in the detector are simulated using GENIE 2.6.2~\cite{Andreopoulos201087}. 
In GENIE, the initial nucleon momentum is selected from distributions in Refs.~\cite{Bodek:1980ar,Bodek:1981wr}.   
Scattering kinematics are calculated in the (off-shell) nucleon rest frame.    
The quasielastic cross section is reduced to account for Pauli blocking.    
For quasielastic and resonance processes, free nucleon form factors are used.  
Quasielastic model details are given in Ref.~\cite{nubarprl}. 
Kinematics for nonresonant inelastic processes are selected from the model of Ref.~\cite{Bodek:2004pc}
which effectively includes target mass and higher twist corrections. 
An empirical correction factor based on charged lepton deep inelastic scattering measurements of 
$F_{2}^{D}/F_2^{\left(n+p\right)}$ and $F_{2}^{Fe}/F_{2}^{D}$ is applied to all structure functions 
as a function of \Xbj, independent of the four-momentum transfer squared \Qsq and \A.  
This accounts for all nuclear effects except those related to neutron excess, which are applied separately.

The \minerva detector's response is simulated by a tuned 
GEANT4-based~\cite{Agostinelli2003250,1610988} simulation.  
The energy scale of the detector is set by ensuring
both detected photon statistics and reconstructed energy deposited by
momentum-analyzed throughgoing muons agree in data and simulation.
Calorimetric constants applied to reconstruct the recoil energy 
are determined by simulation. 
This procedure is cross-checked by comparing data and simulation 
of a scaled-down version of the \minerva detector in a 
low energy hadron test beam~\cite{minerva_nim}.  

